\chapter*{Summary}
\addcontentsline{toc}{chapter}{Summary}

The aim of this thesis was to present, combine, enhance, and apply various edge enhancement algorithms with the purpose of improving the visibility of blood vessels on \gls{dva} images.

First, the topic of spatial domain filtering has been explored. The outcomes were not persuasive: the implementation of gradient filters resulted in an undesired artefact of edge-doubling, whereas the utilisation of adaptive filtering did not produce substantial enhancements neither.

Then, the application of spatial domain filtering showed positive outcomes. The ability to selectively amplify signals within a given band of frequencies facilitated the enhancement of blood vessels with particular widths, while effectively removing noise or larger structures.

The application of \gls{pca} in \gls{dva} appeared to be a potentially advantageous approach for emphasising the vessel region while minimising background signal. However, the outcomes did not demonstrate satisfactory results. By enhancing the domain selection algorithm, it may possible to achieve improved outcomes.

One effective approach for blood vessel segmentation could be the utilisation of neural networks \cite{nasr2018segmentation}. The results generated by such an approach can be simply transferred to a post-processing technique in order to achieve edge enhancing. Another method that was not covered in this thesis is the wavelet transform, which has the ability to decompose signals into localised oscillations in space. This is in contrast to the Fourier transform, which decomposes signals into oscillations that persist across the entire sequence \cite{akram2009blood}.



